\documentclass[12pt,twoside]{amsart}
\usepackage{a4wide, amsmath,amsbsy,amsfonts,amssymb,
stmaryrd,amsthm,mathrsfs,graphicx, amscd, tikz-cd, cancel}
\usepackage[normalem]{ulem}
\usepackage{soul}
\usetikzlibrary{matrix,arrows,decorations.pathmorphing}
\pagestyle{headings}
\markboth{Chapter}{Section}
\usepackage[active]{srcltx}
\usepackage[
	hypertexnames=false,
	hyperindex,
	pagebackref,  
	pdftex,
	breaklinks=true,
	bookmarks=false,
	colorlinks,
	linkcolor=blue,
	citecolor=red,
	urlcolor=red,
]{hyperref}
\newcommand{\opm}[1]{\textcolor{red}{$\star$#1$\star$}}
\newcommand{\change}[1]{\textcolor{blue}{#1}}

\usepackage[all]{xy}
\usepackage{subfig}
\usepackage{tikz}
\usetikzlibrary{shapes,arrows,shadows}
\usetikzlibrary{decorations.markings}
\usepackage{color}
\usepackage[normalem]{ulem}
\usepackage{hyperref}
\usepackage{wrapfig}
\usetikzlibrary{arrows}
\usepackage{pinlabel}
\long\def\symbolfootnote[#1]#2{\begingroup%
\def\thefootnote{\fnsymbol{footnote}}\footnote[#1]{#2}\endgroup}

\renewcommand{\labelenumi}{(\roman{enumi})}


\newtheorem{innercustomthm}{Theorem}
\newenvironment{customtheorem}[1]
  {\renewcommand\theinnercustomthm{#1}\innercustomthm}
  {\endinnercustomthm}
  
  \newtheorem{innercustomcor}{Corollary}
\newenvironment{customcorollary}[1]
  {\renewcommand\theinnercustomcor{#1}\innercustomcor}
  {\endinnercustomcor}




\newtheorem{theorem}{Theorem}[section]
\newtheorem*{NoNumberTheorem}{Theorem}
\newtheorem{proposition}[theorem]{Proposition}
\newtheorem{corollary}[theorem]{Corollary}
\newtheorem{lemma}[theorem]{Lemma}
\newtheorem{claim}[theorem]{Claim}
\newtheorem{sublemma}[theorem]{Sublemma}

\newtheorem{conjecture}[theorem]{Conjecture}
\newtheorem{question}[theorem]{Question}

\newtheorem*{convention}{Convention}
\newtheorem*{notation}{Notation}

\theoremstyle{definition}
\newtheorem{construction}[theorem]{Construction}
\newtheorem{case}[theorem]{Case}

\newtheorem{definition}[theorem]{Definition}
\newtheorem{remark}[theorem]{Remark}
\newtheorem{remarks}[theorem]{Remarks}  
%\newtheorem{notation}[theorem]{Notation}
\newtheorem{example}[theorem]{Example}
\newtheorem{quest}[theorem]{Question}

\newtheorem*{namedtheorem}{\theoremname}
\newcommand{\theoremname}{testing}
\newenvironment{named}[1]{\renewcommand{\theoremname}{#1}
\begin{namedtheorem}}{\end{namedtheorem}}



\setcounter{tocdepth}{1}

%NEWCOMMANDS
\DeclareMathSymbol{\Alpha}{\mathalpha}{operators}{"41}
\DeclareMathSymbol{\Beta}{\mathalpha}{operators}{"42}
\DeclareMathSymbol{\Epsilon}{\mathalpha}{operators}{"45}
\DeclareMathSymbol{\Zeta}{\mathalpha}{operators}{"5A}
\DeclareMathSymbol{\Eta}{\mathalpha}{operators}{"48}
\DeclareMathSymbol{\Iota}{\mathalpha}{operators}{"49}
\DeclareMathSymbol{\Kappa}{\mathalpha}{operators}{"4B}
\DeclareMathSymbol{\Mu}{\mathalpha}{operators}{"4D}
\DeclareMathSymbol{\Nu}{\mathalpha}{operators}{"4E}
\DeclareMathSymbol{\Omicron}{\mathalpha}{operators}{"4F}
\DeclareMathSymbol{\Rho}{\mathalpha}{operators}{"50}
\DeclareMathSymbol{\Tau}{\mathalpha}{operators}{"54}
\DeclareMathSymbol{\Chi}{\mathalpha}{operators}{"58}
\DeclareMathSymbol{\omicron}{\mathord}{letters}{"6F}
\newcommand{\R}{\mathbb{R}}
\newcommand{\C}{\mathbb{C}}
\newcommand{\N}{\mathbb{N}}
\newcommand{\Z}{\mathbb{Z}}
\newcommand{\ZZ}{{\widehat{\mathbb Z}}}
\newcommand{\hZ}{\widehat{\mathbb{Z}}}
\newcommand{\Q}{\mathbb{Q}}
\newcommand{\I}{\mathbb{I}}
\newcommand{\bt}{\bullet}
\newcommand{\Ends}{{\text{Ends}}}
\renewcommand{\tablename}{Diagram}
\newcommand\Res{\operatorname{Res}}
\renewcommand\Im{\operatorname{Im}}
\newcommand\vcd{\operatorname{vcd}}

%OPERATORS

\def\Aut{\operatorname{Aut}}
\def\Sp{\operatorname{Sp}}
\def\GL{\operatorname{GL}}
\def\SL{\operatorname{SL}}
\def\PSL{\operatorname{PSL}}
\def\PGL{\operatorname{PGL}}
\def\Inn{\operatorname{Inn}}
\def\inn{\operatorname{inn}}

\def\Out{\operatorname{Out}}
\def\Map{\operatorname{Map}}
\def\Homeo{\operatorname{Homeo}}
\def\Hom{\operatorname{Hom}}
\def\Der{\operatorname{Der}}
\def\PG{\mathrm{P}\Gamma}
\def\Star{\operatorname{Star}}
\def\Link{\operatorname{Lk}}
\def\Isom{\operatorname{Isom}}
\def\Spec{\operatorname{Spec}}
\def\Fix{\operatorname{Fix}}
\def\P{{\mathbb P}}
\def\H{{\mathbb H}}
\def\M{{\mathbb M}}
\def\k{\Bbbk}
\def\kk{{\overline{\k}}}
\def\cP{{\mathcal P}}
\def\cS{{\mathscr S}}
\def\hS{\hat{\mathscr S}}
\def\cL{{\mathcal L}}
\def\cI{{\mathcal I}}
\def\ccI{{\mathscr I}}
\def\hL{\widehat{\mathcal L}}
\def\cB{{\mathcal B}}
\def\cA{{\mathcal A}}
\def\cQ{{\mathcal Q}}

\def\la{\langle}
\def\ra{\rangle}

\def\cT{{\mathcal T}}
\def\hcT{{\widehat {\mathcal T}}}
\def\bcT{{\overline{\mathcal T}}}
%\def\wT{{\widetilde{\mathcal T}}}
\def\wT{{\widehat{\mathcal T}}}
%\def\dwT{{\partial\widetilde{\mathcal T}}}
\def\tS{\widetilde{S}}
\def\cC{{\mathscr C}}
\def\cF{{\mathcal F}}
\def\hF{{\widehat{\operatorname{F}}}}
%\def\nF{\boldsymbol{\mathcal F}}
\def\nF{\boldsymbol{\mathcal F}}
%\def\wF{{\widetilde{\mathcal F}}}
\def\wF{{\widehat{\mathcal F}}}
\def\cG{{\mathcal G}}
%\def\cH{{\mathcal H}}
\def\cR{{\mathcal R}}
%\def\wH{{\widetilde{\mathcal H}}}
\def\wH{{\widehat{\mathcal H}}}
\def\cM{{\mathcal M}}
\def\cP{{\mathcal P}}
\def\hP{{\widehat\Pi}}
\def\dM{{\partial\mathcal M}}
\def\ccM{{\overline{\mathcal M}}}
%\def\wM{{\widetilde{\mathcal M}}}
\def\wM{{\widehat{\mathcal M}}}
%\def\dwM{{\partial\widetilde{\mathcal M}}}
\def\dwM{{\partial\widehat{\mathcal M}}}
\def\dd{\partial}
\def\a{{\alpha}}
\def\b{{\beta}}
\def\l{{\lambda}}
\def\L{{\Lambda}}
\def\m{{\mu}}

\def\g{{\gamma}}
\def\G{{\Gamma}}
\def\kG{{\check\Gamma}}
\def\kPG{\mathrm{P}{\check\Gamma}}
\def\hG{{\widehat\Gamma}}
\def\hPG{\mathrm{P}{\widehat\Gamma}}
\def\T{{\Theta}}
\def\tt{{\theta}}

\def\d{{\delta}}
\def\s{{\sigma}}
\def\t{{\tau}}
\def\bt{\bullet}

\def\hB{\widehat{B}}
\def\kB{\check{B}}
\def\kC{\check{C}}
\def\hC{\widehat{C}}
\def\II{\mathrm{I}}
\def\hI{{\hat{\mathrm{I}}}}

\def\oC{{\overline C}}
\def\kN{{\check{\Nu}}}

\def\cD{{\mathcal D}}
\def\kcD{{\check{\mathcal D}}}

\def\ssm{\smallsetminus}
\def\ol{\overline}
\def\wh{\widehat}
\def\wt{\widetilde}
\def\ul{\underline}
\def\da{\downarrow}
\def\sr{\stackrel}
\def\hookra{\hookrightarrow}
\def\co{\colon\thinspace}


\begin{document}

\title[A congruence subgroup property for symmetric mapping class groups]{A congruence subgroup property for symmetric mapping class groups}
 
\author{Marco Boggi}
\address{UFF - Instituto de Matem\'atica e Estat\'{\i}stica - Niter\'oi - RJ 24210-200; Brazil}
\email{marco.boggi@gmail.com}


\begin{abstract}    
We prove the congruence subgroup property for the centralizer of a finite subgroup $G$ in the mapping class group of a hyperbolic oriented 
and connected surface of finite topological type $S$ such that the genus of the quotient surface $S/G$ is at most $2$. 
As an application, we show that torsion elements in the mapping class group of a surface of genus $\leq 2$ are conjugacy distinguished.
\bigskip

\noindent {\bf AMS Math Classification:} 57K20, 20F65, 14G32.

\end{abstract}

\maketitle


\section{Introduction}\label{Intro}
\subsection{The congruence subgroup property}
Let $S$ be an oriented connected surface of negative Euler characteristic and finite type and let $G$ be a finite group acting faithfully on $S$. 
The group $G$ then identifies with a finite subgroup of the mapping class group $\G(S)$ of the surface $S$ and let us denote by $\G(S)^G$ 
the centralizer of $G$ in $\G(S)$. Let $\Pi$ be the fundamental group of $S$ for some choice of base point. 
There is a natural faithful representation $\rho_S\co\G(S)\hookra\Out(\Pi)$ 
which restricts to a  faithful representation $\rho_{S,G}\co\G(S)^G\hookra\Out(\Pi)$. 

For an abstract group $L$, let us denote by $\wh{L}$ its profinite completion. The homomorphisms $\rho_S$ and $\rho_{S,G}$ then induce the
homomorphisms of profinite groups:
\begin{equation}\label{proGrepr}
\hat\rho_S\co\hG(S)\to\Out(\hP)\hspace{1cm}\mbox{and}\hspace{1cm}\hat\rho_{S,G}\co\hG(S)^G\to\Out(\hP).
\end{equation}
The \emph{congruence subgroup problem for the mapping class group $\G(S)$ (resp.\ for $\G(S)^G$)} 
asks whether the homomorphism $\hat\rho_S$ (resp.\ $\hat\rho_{S,G}$) is faithful. 

More conventionally, this question can be reformulated by asking whether, for every finite index subgroup 
$\G$ of $\G(S)$ (resp.\ of $\G(S)^G$), there exists a finite index characteristic subgroup $K$ of $\Pi$ such that 
the kernel of the induced representation $\G(S)\to\Out(\Pi/K)$ (resp.\ $\G(S)^G\to\Out(\Pi/K)$) is contained in $\G$.
The first question is known to have a positive
answer for $g(S)\leq 2$ (cf.\ \cite{Asada}, \cite{hyp}, \cite{CongTop}), where $g(S)$ denotes the genus of $S$. 

Let $S_{/G}:=S/G$ be the quotient of the surface $S$ for the action of $G$ and let $p_G\co S\to S_{/G}$ be the quotient map. 
Let $R$ and $B$ be, respectively, the ramification and the branch locus of $p_G$ and put $S^\circ:=S\ssm R$ and $S_{/G}^\circ:=S_{/G}\ssm B$.
The main result of the paper is then (cf.\ Theorem~\ref{congrel} for a slightly more general result):

\begin{customtheorem}{A}Let $S$ be an oriented connected hyperbolic surface of finite type and $G$ a finite group acting faithfully on $S$.
If the congruence subgroup property holds for the mapping class group $\G(S_{/G}^\circ)$ of the surface $S_{/G}^\circ$, then it holds for 
the centralizer $\G(S)^G$ of $G$ in the mapping class group $\G(S)$ of $S$.
\end{customtheorem}

In particular, we have:

\begin{customcorollary}{B}If $g(S_{/G})\leq 2$, then the congruence subgroup property holds for $\G(S)^G$.
\end{customcorollary}


\subsection{Conjugacy separability of torsion elements in mapping class groups}
An element $x$ of a group $G$ is \emph{conjugacy distinguished} if, for any $y\in G$, which is not conjugated to $x$ in $G$ 
(briefly $y\not\sim_G x$), there is a finite quotient $\pi\co G\to L$ such that $\pi(y)\not\sim_L \pi(x)$. 

Similarly, a subgroup $H$ of $G$ is \emph{subgroup conjugacy distinguished} if, for any subgroup $K$ of $G$, 
which is not conjugated to $H$ in $G$ (briefly $H\not\sim_G K$), there is a finite quotient $\pi\co G\to L$ such that $\pi(H)\not\sim_L \pi(K)$.

For a subgroup $H$ of a group $G$, we denote by $Z_H(G)$ and $N_H(G)$, respectively, the centralizer and the normalizer of $H$ in $G$. 
For a subset $A$ of a profinite group $R$, we denote by $\ol{A}$ the closure of $A$ for the profinite topology of $R$. We then have
(cf.\ Theorem~\ref{torsionrelative} for a slightly more general result):

\begin{customtheorem}{C}Let $S$ be an oriented connected hyperbolic surface of finite type of genus $\leq 2$. %Then:
\begin{enumerate}
\item The natural monomorphism $\G(S)\hookra\hG(S)$ induces a bijection between conjugacy classes
of solvable finite subgroups. In particular, solvable finite subgroups of $\G(S)$ are subgroup conjugacy distinguished. 
\item Torsion elements of $\G(S)$ are conjugacy distinguished.
\item For every finite subgroup $G$ of $\G(S)$, there holds $Z_{\hG(S)}(G)=\ol{Z_{\G(S)}(G)}$.\\ 
%If, moreover, $G$ is solvable, we also have $N_{\hG(S,\dd S)}(G)=\ol{N_{\G(S,\dd S)}(G)}$.
\end{enumerate}
\end{customtheorem}

 


\section{Preliminary results}
\subsection{Normalizers and centralizers of finite subgroups of the mapping class group of a surface}\label{centrasec}
It is clearly not restrictive to assume that $S$ has empty boundary and then that it is obtained from a closed surface removing a finite set of points.
With the notations of Section~\ref{Intro}, the cover $p_G\co S\to S_{/G}$ determines a normal finite index subgroup $K\cong\pi_1(S^\circ)$ 
of the fundamental group $\Pi_{/G}: =\pi_1(S_{/G}^\circ)$ (we omit base points). 

Let $\G(S_{/G}^\circ)_K$ be the finite index subgroup of the mapping class group $\G(S_{/G}^\circ)$ which consists of the mapping classes which 
preserve the normal subgroup $K$ of $\Pi_{/G}$. The normalizer $N_{\G(S)}(G)$ of the finite subgroup $G\cong\left.\Pi_{/G}\right/K$ 
in the mapping class group $\G(S)$ is then described by the short exact sequence:
\begin{equation}\label{normalizerG}
1\to G\to N_{\G(S)}(G)\to\G(S_{/G}^\circ)_K\to 1.
\end{equation}

Let $Z(G)$ be the center of the group $G$ and let $\G[G]$ be the kernel of the natural representation $\G(S_{/G}^\circ)_K\to\Out(G)$.
The centralizer $\G(S)^G\leq N_{\G(S)}(G)$ of the subgroup $G$ in the mapping class group $\G(S)$ 
is then described by the short exact sequence:
\begin{equation}\label{centralizerG}
1\to Z(G)\to\G(S)^G\to\G[G]\to 1.
\end{equation}

The group $G$ is the covering transformation group of the unramified cover $p_G^\circ\co S^\circ\to S_{/G}^\circ$ and thus also identifies 
with a finite subgroup of the mapping class group $\G(S^\circ)$ of $S^\circ$.
From the short exact sequences~\eqref{normalizerG} and~\eqref{centralizerG}, it follows:

\begin{proposition}\label{comparison}The natural epimorphism of mapping class groups $q^\circ\co\G(S^\circ)\to\G(S)$, associated to the embedding
of $S^\circ$ in $S$, induces the isomorphisms: 
\[N_{\G(S^\circ)}(G)\cong N_{\G(S)}(G)\hspace{1cm}\mbox{and}\hspace{1cm}\G(S^\circ)^G\cong\G(S)^G.\]
\end{proposition}

\begin{proof}Let us just prove the first isomorphism, the second one is proved similarly. By the short exact sequence~\eqref{normalizerG},
applied to the surfaces $S^\circ$ and $S$, respectively, there is a commutative diagram with exact rows:
\[\begin{array}{ccccccc}
1\to&G&\to&N_{\G(S^\circ)}(G)&\to&\G(S_{/G}^\circ)_K&\to 1\,\\
&\parallel&&\da^{q^\circ}&&\parallel&\\
1\to&G&\to&N_{\G(S)}(G)&\to&\G(S_{/G}^\circ)_K&\to 1,
\end{array}\]
from which the conclusion follows.
\end{proof}



\subsection{The congruence topology}\label{congtopology}
The \emph{congruence topology} on $\G(S)$ is the subspace topology induced by the standard profinite topology on $\Out(\hP)$ 
via the natural monomorphism $\G(S)\hookra\Out(\hP)$. The congruence topology on $\G(S)$ then induces a topology on its subgroups
which we also call the \emph{congruence topology}.

Let $\G(m)^G$ be the kernel of the natural representation $\G(S)^G\to\GL(H_1(S,\Z/m))$ and let $\G[G,m]$
be its image via the natural epimorphism $\G(S)^G\to\G[G]$. For $m\geq 3$, it is well known that the action of $G$ on $H_1(S,\Z/m)$
is faithful. Hence, there holds $\G(m)^G\cap G=\{1\}$ and so the map
$\G(S)^G\to\G[G]$ induces an isomorphism $\G(m)^G\cong\G[G,m]$. 

\begin{lemma}\label{congsub}For all $m\geq 2$, the finite index subgroups $\G(m)^G$ of $\G(S)^G$ and $\G[G,m]$ of $\G(S_{/G}^\circ)$
are both open for the respective congruence topologies.
\end{lemma}

\begin{proof}The statement is obvious for $\G(m)^G$. For $\G[G,m]$,
it can be proved with the same argument of \cite[Corollary~2.3]{sym}. 
\end{proof}


\subsection{The centralizer of a finite subgroup in the profinite completion of a good group}
Let $\iota\co G\to\wh{G}$ be the natural homomorphism of a group to its profinite completion. The group $G$ is \emph{good}
if, for every finite discrete $\wh{G}$-module $M$, the homomorphism induced on cohomology $\iota^\ast\co H^i(\wh{G},M)\to H^i(G,M)$ 
is an isomorphism for all $i\geq 0$. Let us recall the following definition (cf.\ \cite[Definition~1.2]{BZ}):

\begin{definition}\label{fct}A group $G$ is of \emph{finite virtual cohomological type} if:
\begin{enumerate}
\item $G$ has finite virtual cohomological dimension;
\item for every finite $G$-module $M$, the cohomology $H^i(G,M)$ is finite for $i\geq 0$.
\end{enumerate}
\end{definition}

As a corollary of \cite[Theorem~A]{BZ}, we then have:

\begin{proposition}\label{centgood}Let $G$ be a good, residually finite group of finite virtual cohomological type 
and let $\iota\co G\hookra\wh{G}$ be the natural monomorphism. Then, for every finite subgroup $H$ of $G$, 
we have the identity $Z_{\wh{G}}(\iota(H))=\ol{\iota(Z_G(H))}$.
\end{proposition}

\begin{proof}By \cite[(iii) of Theorem~A]{BZ}, for all primes $p>0$ and every $p$-torsion element $x$ of $G$, there holds 
$Z_{\wh{G}}(\iota(x))=\ol{\iota(Z_G(x))}$. The finite group $H$ is generated by a finite subset $\{x_1,\dots,x_k\}$ of $p$-torsion elements 
(for, possibly, different primes $p>0$). We then have:
\[Z_{\wh{G}}(\iota(H))=\bigcap_{i=1}^kZ_{\wh{G}}(\iota(x_i))=\bigcap_{i=1}^k\ol{\iota(Z_G(x_i))}=\ol{\iota(Z_G(H))}.\]
\end{proof}


\section{Proof of Theorem~A}
\subsection{The case of a free action}\label{freeaction}
The first case to consider for the proof of Theorem~A is the one in which $G$ acts freely on $S$, that is to say $R=B=\emptyset$ and so
$\Pi=K$ identifies with a normal finite index subgroup of $\Pi_{/G}=\pi_1(S_{/G})$. 

The congruence subgroup property for $\G(S)^G$ is equivalent to the statement that the congruence topology on this group, as defined 
in Section~\ref{congtopology}, is equivalent to the profinite topology. Since, by Lemma~\ref{congsub}, the subgroup $\G(m)^G$ of $\G(S)^G$ 
is open for the congruence topology, the congruence subgroup property for $\G(S)^G$ follows if we show that the homomorphism:
\[\hat\rho_{S,G}^{(m)}\co\hG(m)^G\to\Out(\hP),\]
induced by the restriction of $\rho_{S,G}$ to $\G(m)^G$, is injective for some $m\geq 2$. 

In order to prove the above claim, we will need the well known lemma: 

\begin{lemma}\label{restriction}Let $\hP$ be a nonabelian profinite surface group, $N$ an open normal subgroup of $\hP$ and $f\in\Aut(\hP)$.
If $f$ restricts to the identity on $N$, then $f$ is the identity.
\end{lemma}

\begin{proof}The proof is essentially the same of \cite[Lemma~8] {Asada}.
Let $a_1,a_2\in\hP$ be distinct simple generators. Then, for some $n\in\N^+$, we have that $a_1^n,a_2^n\in N$ and,
since $N$ is normal, also that $x a_i^n x^{-1}\in N$, for all $x\in\hP$ and $i=1,2$. 

By hypothesis, we have that $f(x a_i^n x^{-1})=f(x)a_i^n f(x)^{-1}=x^{-1}a_i^n x$ and then $f(x)x^{-1}\in N_\hP(\la a_i^n\ra)$, for $i=1,2$. 
Since, by \cite[Proposition~3.6]{Magnus}, there holds $N_\hP(\la a_i^n\ra)=\la a_i\ra$, for $i=1,2$, and $\la a_1\ra\cap\la a_2\ra=\{1\}$,
the conclusion follows.
\end{proof}

We will also need the following remarks.
Since, by hypothesis, the representation $\hat\rho_{S_{/G}}\co\hG(S_{/G})\to\Out(\hP_{/G})$ is faithful, its restriction 
$\hat\rho_{S_{/G}}^{(m)}\co\hG[K,m]\to\Out(\hP_{/G})$ is faithful as well for all $m\geq 2$.
For $m\geq 3$, the natural isomorphism $\G(m)^G\cong\G[G,m]$ induces an isomorphism of profinite completions 
$\hG(m)^G\cong\hG[K,m]$ and $\hP\cong\wh{K}$ identifies with a normal finite index subgroup of 
$\hP_{/G}$ in a way that the representations $\hat\rho_{S,G}^{(m)}$ and $\hat\rho_{S_{/G}}^{(m)}$ are compatible via the above isomorphisms.

Let then $x\in\ker\hat\rho_{S,G}^{(m)}$, put $y:=\hat\rho_{S_{/G}}^{(m)}(x)\in\Out(\hP_{/G})$ and let $\tilde y\in\Aut(\hP_{/G})$ be a lift.
Since $x\in\ker\hat\rho_{S,G}^{(m)}$, by the remarks above, after possibly composing with an inner automorphism of $\hP_{/G}$, 
we can assume that the restriction of $\tilde y$ to $\wh{K}\equiv\hP$ is the identity. By Lemma~\ref{restriction}, this implies that $\tilde y$ 
is the identity and so that $y=1$. Since the representation $\hat\rho_{S_{/G}}^{(m)}$ is faithful, it follows that $x=1$ as well. 
This concludes the proof of Theorem~A for the case of a free action.


\subsection{The general case}
The completions of the groups $\G(S)$ and $\G(S)^G$ with respect to the congruence topologies which we introduced in Section~\ref{congtopology}
are called \emph{procongruence completions} and denoted by $\kG(S)$ and $\kG(S)^G$, respectively. The group $\kG(S)$ is also called the
\emph{procongruence mapping class group}. The congruence subgroup problems for $\G(S)$ and $\G(S)^G$ can then be formulated as the
questions of whether the natural epimorphisms  $\hG(S)\to\kG(S)$ and $\hG(S)^G\to\kG(S)^G$ are, respectively, isomorphisms.

By the results of Section~\ref{freeaction} and with the notations of Section~\ref{Intro}, we know that there is a natural isomorphism 
$\hG(S^\circ)^G\cong\kG(S^\circ)^G$. Therefore, in order to prove the general case of Theorem~A, it is enough to show that the 
natural isomorphism $q^\circ\co\G(S^\circ)^G\sr{\sim}{\to}\G(S)^G$ (cf.\ Proposition~\ref{comparison})
induces an isomorphism $\check{q}^\circ\co\kG(S^\circ)^G\sr{\sim}{\to}\kG(S)^G$ between the respective procongruence completions.

The set $R$ of ramification points of the quotient map $p_G\co S\to S_{/G}$ determines a point, which we denote by $[R]$, 
on the unordered configuration space $S[r]$ of $r$ distinct points on the surface $S$, where $r:=\sharp R$. 

Let $\Pi[r]:=\pi_1(S[r],[R])$. From \cite[Corollary~4.7 and Proposition~4.9]{CongTop}, 
it then follows that there is a short exact sequence of profinite groups:
\begin{equation}\label{Birman}
1\to\hP[r]\to\kG(S^\circ)\sr{\check{q}^\circ}{\to}\kG(S)\to 1.
\end{equation}

In order to prove Theorem~A, it is then enough to show that the subgroup $\kG(S^\circ)^G$ of $\kG(S^\circ)$ has trivial intersection with
the image of $\hP[r]$ in $\kG(S^\circ)$. Since $\kG(S^\circ)^G$ is contained in the centralizer $Z_{\kG(S^\circ)}(G)$, 
the conclusion follows if we prove that $Z_{\kG(S^\circ)}(G)\cap\hP[r]=\{1\}$.

By Proposition~\ref{comparison} and the short exact sequence~\eqref{Birman}, the subgroup $G$ of $\G(S^\circ)$ has trivial intersection  
with $\hP[r]$, so that $\hP[r]\cdot G\cong\hP[r]\rtimes G$ identifies with the profinite completion of $\Pi[r]\cdot G\cong\Pi[r]\rtimes G$.
We are then reduced to prove the following lemma:

\begin{lemma}\label{Gcentconf}There holds $Z_{\hP[r]\cdot G}(G)=Z(G)$.
\end{lemma}

\begin{proof}Proposition~\ref{comparison} implies that $\G(S^\circ)^G\cap\Pi[r]=\{1\}$ and then that $Z_{\Pi[r]\cdot G}(G)=Z(G)$.
Since $\Pi[r]\cdot G$ contains, as a finite index subgroup, the fundamental group of the \emph{ordered} configuration space of 
$n$ points on the surface $S$, which is a good, residually finite group of finite virtual cohomological type, by the Lyndon-Hochschild-Serre 
spectral sequence, it follows that $\Pi[r]\cdot G$ has the same properties. Hence, we conclude by Proposition~\ref{centgood}.
\end{proof}




\section{The congruence subgroup property for relative mapping class groups}
For $S$ a surface as above, let us denote by $\dd S$ its boundary and by $\G(S,\dd S)$ the \emph{relative} mapping class group of the 
surface with boundary $(S,\dd S)$, that is to say the group of relative, with respect to the boundary $\dd S$, isotopy classes of orientation 
preserving diffeomorphisms of $S$. Of course, if $\dd S=\emptyset$, then $\G(S,\dd S)=\G(S)$.

Let us observe that, for $\dd S\neq\emptyset$, the natural representation $\G(S,\dd S)\to\Out(\Pi)$ is not faithful. 
Hence, the congruence subgroup property, as formulated in Section~\ref{Intro} does not make sense for relative mapping class groups. 
For this reason, in \cite[Section~2.7]{Aut}, we introduced the following construction.

Let $\s:=\{\delta_1,\ldots,\delta_k\}$ be the set of connected components of $\dd S$ and let $\tS$ be the surface (without boundary) 
obtained from $S$ attaching to each curve $\delta_i$ a  a $2$-punctured closed disc $D_i$, for $i=1,\dots k$. Let then $\G(\tS)_\s$ 
be the stabilizer in $\G(\tS)$ of the set of isotopy classes of the curves contained in $\s$. 

The relative mapping class group $\G(S,\dd S)$ and the product of relative mapping class groups $\prod_{i=1}^k\G(D_i,\dd D_i)$ then identify,
respectively, with a subgroup and a normal subgroup of the stabilizer $\G(\tS)_\s$ and there holds:
\begin{equation}\label{embedding}
\G(\tS)_\s=(\prod_{i=1}^k\G(D_i,\dd D_i))\cdot\G(S,\dd S).
\end{equation}

The identity~\eqref{embedding} identifies the relative mapping class group $\G(S,\dd S)$ with a subgroup of the mapping class group $\G(\tS)$ and we
define the \emph{congruence topology on $\G(S,\dd S)$} as the topology induced, via this embedding, by the congruence topology of $\G(\tS)$. 
The \emph{relative procongruence mapping class group} $\kG(S,\dd S)$ is the associated completion of $\G(S,\dd S)$.

The \emph{congruence subgroup problem} for $\G(S,\dd S)$ then asks whether the natural epimorphism $\hG(S,\dd S)\to\kG(S,\dd S)$ is an isomorphism. 

For a finite group acting faithfully on $S$, we let $\G(S,\dd S)^G$ be the centralizer of $G$ in $\G(S,\dd S)$. We define the congruence
topology on $\G(S,\dd S)^G$ as the one induced by the congruence topology of $\G(S,\dd S)$ and we denote by $\kG(S,\dd S)^G$ the
associated completion. We can then formulate the corresponding congruence subgroup problem for $\G(S,\dd S)^G$.

Let us observe that there is a natural short exact sequence:
\begin{equation}\label{relexseq1}
1\to\prod_{i=1}^k\tau_{\d_i}^\Z\to\G(S,\dd S)\to\G(S)\to 1,
\end{equation}
where $\tau_{\d_i}^\Z\cong\Z$ is the cyclic subgroup of $\G(S,\dd S)$ topologically generated by the Dehn twist about the curve $\d_i$, for $i=1,\dots k$.
In \cite[Section~2.7]{Aut}, it was then observed that the procongruence completion of $\G(S,\dd S)$ induces the short exact sequence:
\begin{equation}\label{relexseq2}
1\to\prod_{i=1}^k\tau_{\d_i}^\ZZ\to\kG(S,\dd S)\to\kG(S)\to 1,
\end{equation}
where $\tau_{\d_i}^\ZZ\cong\ZZ$ is the closure of $\tau_{\d_i}^\Z$ in the procongruence completion $\kG(S,\dd S)$, for $i=1,\dots k$. 

We then have the following generalization of Theorem~A and Corollary~B:

\begin{theorem}\label{congrel}Let $S$ be an oriented connected hyperbolic surface of finite type and $G$ a finite group acting faithfully on $S$.
If the congruence subgroup property holds for the mapping class group $\G(S_{/G}^\circ)$ of the surface $S_{/G}^\circ$, 
then it also holds for the centralizer $\G(S,\dd S)^G$ of $G$ in the relative mapping class group $\G(S,\dd S)$. 
In particular, for $g(S_{/G})\leq 2$, there is a natural isomorphism $\hG(S,\dd S)^G\cong\kG(S,\dd S)^G$.
\end{theorem}

\begin{proof}The group $G$ acts by conjugation on the normal abelian subgroup $\prod_{i=1}^k\tau_{\d_i}^\Z$ of $\G(S,\dd S)$.
Let us denote by $A^G$ the invariants for the action of a group $G$ on an abelian group $A$. We then have the short exact sequence:
\[1\to(\prod_{i=1}^k\tau_{\d_i}^\Z)^G\to\G(S,\dd S)^G\to\G(S)^G\to 1.\]
Similarly, for the procongruence completion, we get the short exact sequence:
\[1\to(\prod_{i=1}^k\tau_{\d_i}^\ZZ)^G\to\kG(S,\dd S)^G\to\kG(S)^G\to 1,\]
where the subgroup of invariants $(\prod_{i=1}^k\tau_{\d_i}^\ZZ)^G$ identifies with the closure of the subgroup $(\prod_{i=1}^k\tau_{\d_i}^\Z)^G$
in $\kG(S,\dd S)^G$. Theorem~\ref{congrel} then follows from the two above exact sequences and Theorem~A if we show that the profinite group 
$(\prod_{i=1}^k\tau_{\d_i}^\ZZ)^G$ is the profinite completion of the group $(\prod_{i=1}^k\tau_{\d_i}^\Z)^G$. But this is clear 
since the profinite completion of a finitely generated abelian group induces the profinite completion on each of its subgroups.
\end{proof}



\section{Conjugacy separability}
We will prove the following more general version of Theorem~C:

\begin{theorem}\label{torsionrelative}For $S$ an oriented connected hyperbolic surface of finite type of genus $\leq 2$, we have:
\begin{enumerate}
\item The natural monomorphism $\G(S,\dd S)\hookra\hG(S,\dd S)$ induces a bijection between conjugacy classes
of solvable finite subgroups. In particular, solvable finite subgroups of $\G(S,\dd S)$ are subgroup conjugacy distinguished. 
\item Torsion elements of $\G(S,\dd S)$ are conjugacy distinguished.
\item For every finite subgroup $G$ of $\G(S,\dd S)$, there holds $Z_{\hG(S,\dd S)}(G)=\ol{Z_{\G(S,\dd S)}(G)}$.\\ 
\end{enumerate}
\end{theorem}

\begin{proof}By \cite[Corollary~C]{BZ}, it is enough to prove that, for $g(S)\leq 2$, the relative mapping class group $\G(S,\dd S)$ 
is good and that, for every finite index subgroup $G$ of $\G(S,\dd S)$, the following conditions are satisfied:
\begin{enumerate}
\item the normalizer $N_{\G(S,\dd S)}(G)$ is also a good group;
\item the closure of the normalizer $N_{\G(S,\dd S)}(G)$ in $\hG(S,\dd S)$ identifies with the profinite completion of $N_{\G(S,\dd S)}(G)$.
\end{enumerate}

The claim that $\G(S,\dd S)$ is good for $g(S)\leq 2$ follows from the fact that $\G(S)$ is a good group for $g(S)\leq 2$ and from the 
Lyndon-Hochschild-Serre spectral sequence applied to the natural homomorphism from the short exact sequence~\eqref{relexseq1} 
to the short exact sequence~\eqref{relexseq2}.

Let us then prove that the normalizer $N_{\G(S,\dd S)}(G)$ is also a good group for $g(S_{/G})\leq 2$ (and so, in particular, for $g(S)\leq 2$).
Let us observe first that there is the following version of the short exact sequence~\eqref{normalizerG} for relative mapping class groups:
\begin{equation}\label{relnorm}
1\to G\to N_{\G(S,\dd S)}(G)\to\G(S_{/G}^\circ,\dd S_{/G}^\circ)_K\to 1.
\end{equation}

Since $\G(S_{/G}^\circ,\dd S_{/G}^\circ)_K$ is a finite index subgroup of $\G(S_{/G}^\circ,\dd S_{/G}^\circ)$ and the latter, as we proved above,
is good for $g(S_{/G})\leq 2$, from Shapiro's lemma, it follows that $\G(S_{/G}^\circ,\dd S_{/G}^\circ)_K$ is also good for $g(S_{/G})\leq 2$.
By the Lyndon-Hochschild-Serre spectral sequence applied to the natural homomorphism which maps the short exact sequence~\eqref{relnorm} 
to its profinite completion, we then have that $N_{\G(S,\dd S)}(G)$ is a good group for $g(S_{/G})\leq 2$. 

By Theorem~\ref{congrel}, for $g(S_{/G})\leq 2$, the closure of the centralizer $\G(S,\dd S)^G$ in $\hG(S,\dd S)$ 
identifies with the profinite completion of $\G(S,\dd S)^G$. Since  $\G(S,\dd S)^G$ is a finite index subgroup of 
the normalizer $N_{\G(S,\dd S)}(G)$, it follows that, for $g(S_{/G})\leq 2$ (and so, in particular, for $g(S)\leq 2$), 
the closure of the normalizer $N_{\G(S,\dd S)}(G)$ in $\hG(S,\dd S)$ identifies with the profinite completion of $N_{\G(S,\dd S)}(G)$.
\end{proof}
 \bigskip
 
%\subsection*{Conflict of interest declaration:} The author states that there is no conflict of interest.


\begin{thebibliography}{99}

%\vskip .1cm
%\bibitem {[A]} M.P.Anderson. Exactness properties of profinite completion functors.
%Topology {\bf 13} (1974), 229-239.

%\bibitem{ACV}D. Abramovich, A. Corti and A. Vistoli. \textsl{Twisted bundles and admissible covers}. 
%Comm. Algebra n. {\bf 31} (2003), 3547--3618.

%\bibitem{AGV}
%D. Abramovich, T. Graber and A. Vistoli, {\em Gromov-Witten theory of Deligne-Mumford stacks}, 
%American J. Math. 130 (2008), 1337--1398. 

%\bibitem{AV} D. Abramovich, A. Vistoli. \textsl{Compactifying the space of stable maps.} J. Amer. Math. Soc., vol. {\bf 15} (2001), n. 1, 27--75.

\bibitem{Asada}M. Asada. \textsl{The faithfulness of the monodromy representations associated
with certain families of algebraic curves.} J. Pure Appl. Algebra {\bf 159}, (2001), 123--147. 

%\bibitem{BP}M. Boggi, M. Pikaart. \textsl{Galois covers of moduli of curves}. Comp. Math. {\bf 120} (2000), 171-191.

%\bibitem{ProTeich} M. Boggi. \textsl{Profinite Teichm{\"u}ller theory.} Math. Nachr. {\bf 279} (2006), 953--987.

%\bibitem{BP} M. Boggi, P. Lochak. \textsl{Profinite complexes of curves, their automorphisms and anabelian properties of moduli stacks of curves.} 
%Preprint.

%\bibitem{mod}M.~Boggi. \textsl{Fundamental groups of moduli stacks of stable curves of compact type.}
%Geometry \& Topology {\bf 13} (2009), 247--276. 

\bibitem{hyp} M. Boggi. \textsl{The congruence subgroup property for the hyperelliptic 
Teichm\"uller modular group: the open surface case}. Hiroshima Math. J. {\bf 39} (2009), 351--362.  

%\bibitem{ProCong} M. Boggi. \textsl{On the procongruence completion of the Teichm\"uller modular groups.}
%Trans. Amer. Math. Soc. {\bf 366} (2014), 5185--5221. 

\bibitem{sym} M. Boggi. \textsl{Galois coverings of moduli spaces of curves and loci of curves with symmetry.} 
Geom. Dedicata {\bf 168} (2014), 113--142. 

%\bibitem{[BZ1]} M. Boggi, P. Zalesskii. \textsl{Characterizing closed curves on Riemann surfaces via 
%homology groups of coverings.}  Int. Math. Res. Not. IMRN 2017, 944--965.

\bibitem{Magnus} M. Boggi, P.A. Zalesskii. \textsl{A restricted Magnus property for profinite surface groups.}
Trans. Amer. Math. Soc. {\bf 371} (2019), 729--753.

\bibitem{BZ} M.~Boggi, P.A.~Zalesskii. \textsl{Finite subgroups of the profinite completion of good groups}. 
Bull. London Math. Soc. vol {\bf 57}, no. 1 (2025), 236--255.

\bibitem{CongTop}M. Boggi. \textsl{Congruence topologies on the mapping class groups}. J. Algebra {\bf 546} (2020), 518--552. 

%\bibitem{HyperNotes}M. Boggi. \textsl{Notes on hyperelliptic mapping class groups.} 
%\url{https://arxiv.org/abs/2110.13534v3} (2023).

%\bibitem{BF}M.~Boggi. L.~Funar.  \textsl{Automorphisms of procongruence curve and pants complexes}. J. Topol. (2023): 
%\url{http://dx.doi.org/10.1112/topo.12306}.

\bibitem{Aut}M. Boggi. \textsl{Automorphisms of profinite mapping class groups}. \url{https://arxiv.org/abs/2011.15075} (2024).

%\bibitem{Collas}
%B.~Collas. \textsl{Action of the Grothendieck-Teichm\"uller group on torsion elements of full Teichm\"uller modular groups in genus zero.} 
%Journal de th\'eorie des nombres de Bordeaux, Volume {\bf 24} (2012) no. 3,  605--622. 

%\bibitem{Drinfeld} V.G. Drinfeld. \textsl{On quasitriangular quasi-Hopf algebras and a group closely connected with $\mathrm{Gal}(\ol{\Q}/\Q)$.}
%Leningrad Math. J. Vol. {\bf 2} (1991), No. 4, 829--860.

%\bibitem{EK} C.J. Earle, I. Kra. \textsl{On isometries between Teichm\"uller spaces.}
%Duke Math. J. {\bf 41} (1974), 583--591.

%\bibitem{FM}B. Farb and D. Margalit. \textsl{A primer on mapping class groups}. Princeton Math. Series, Vol. {\bf 49}, 2012. 

%\bibitem{[FLP]} A.Fathi, F.Laudenbach, V.Po{\'e}naru eds.{\it Travaux de Thurston sur les surfaces} (seconde {\'e}dition), 
%Ast{\'e}risque {\bf 66-67}, SMF Publ., 1991.

%\bibitem{Gervais} S. Gervais. \textsl{A finite presentation of the mapping class group of a punctured surface.}
%Topology {\bf 40} (2001), 703--725.

%\bibitem{Esquisse}A. Grothendieck. \textsl{Esquisse d'un Programme.} In {\it Geometric Galois Actions I}, 
%edited by L. Schneps and P. Lochak, London Math. Soc. LNS {\bf 242} (1997), 5--48.

%\bibitem{[G]}A. Grothendieck. \textsl{Brief an G. Faltings}. In {\it Geometric Galois Actions I}, 
%edited by L. Schneps and P. Lochak, London Math. Soc. LNS {\bf 242} (1997), 49--58.

%\bibitem{[GKM]} A. Gibney, S. Keel, I. Morrison. \textsl{Towards the ample cone of $\ccM_{g,n}$.}
%J. Amer. Math. Soc.  {\bf 15} (2002), 273--294.

%\bibitem{[H1]} J.L.Harer. \textsl{Stability of the homology of the mapping class groups of orientable 
%surfaces.} Annals of Mathematics {\bf 121} (1985), 215-249.

%\bibitem{[H2]} J.L.Harer. \textsl{The cohomology of the moduli spaces of curves.}
%In {\it Theory of Moduli}, LN {\bf 1337} (1988), Springer, 138-221.

%\bibitem{HS} D. Harbater, L. Schneps. \textsl{Fundamental groups of moduli and the Grothendieck-Teichm\"uller group.} 
%Trans. Amer. Math. Soc. {\bf 352} (2000), 3117--3148.

%\bibitem{Harvey} W.J. Harvey. \textsl{Boundary structure of the modular group}. In \emph{Riemann surfaces and related topics:
%Proceedings of the 1978 Stony Brook Conference}. (State Univ. New York, Stony Brook, N.Y., 1978),
%Ann. of Math. Stud. {\bf 97}, Princeton Univ. Press, Princeton, N.J. (1981), 245--251.

%\bibitem{HLS} A. Hatcher, P. Lochak, L. Schneps. \textsl{On the Teichm{\"u}ller tower 
%of mapping class groups.} J. Reine Angew. Math. {\bf 521} (2000), 1--24.

%\bibitem{HM}Y.\ Hoshi, S.\ Mochizuki. \textsl{On the combinatorial anabelian geometry
%of nodally nondegenerate outer representations}. Hiroshima Math. J.
%Volume {\bf 41}, Number 3 (2011), 275--342.

%\bibitem{HMM}Y.\ Hoshi, A. Minamide, S.\ Mochizuki. \textsl{Group-theoreticity of numerical invariants
%and distinguished subgroups of configuration space groups}. Kodai Math. J. {\bf 45} (2022), 295--348.

%\bibitem{[HM]} Y.Hoshi, S.Mochizuki. Topics surrounding the combinatorial anabelian 
%geometry of hyperbolic curves I: inertia groups and profinite Dehn twists. 
%Advanced Studies in Pure Mathematics {\bf 63} (2012), 659-811.

%\bibitem{Huisman} J. Huisman. \textsl{The equivariant fundamental group, uniformization of real algebraic curves, and global complex analytic
%coordinates on Teichm\"uller spaces.} Ann. Fac. Sci. Toulouse Math. $6^e$ s\'erie, tome {\bf 10}, $n^o 4$ (2001), 659--682.

%\bibitem{IN} Y. Ihara, H. Nakamura. \textsl{Some illustrative examples for anabelian geometry in high dimensions.} 
%In {\it Geometric Galois Actions I}, edited by L. Schneps and P. Lochak, Lon. Math. Soc. LNS {\bf 242} (1997), 127-138.

%\bibitem{Ihara}Y. Ihara. \textsl{Some details on the $\GT$-action on $\widehat{B}_n$}. Appendix to: Y. Ihara and
%M. Matsumoto. \textsl{On Galois Actions on Profinite Completions of Braid Groups}. In
%{\it Recent Developments in the Inverse Galois Problem}. Contemp. Math. {\bf 186}, AMS (1995), 173--200.

%\bibitem{[I1]} N.V. Ivanov. \textsl{Subgroups of Teichm{\"u}ller modular groups}.
%Translations of Mathematical Monographs {\bf 115}, AMS Publ., 1992.

%\bibitem{[I2]} N.V. Ivanov. \textsl{Automorphisms of complexes of curves and of Teichm{\"u}ller spaces.}
%Int. Math. Res. Not. IMRN {\bf 14} (1997), 651--666.

%\bibitem{Jarden}  M. Jarden. \textsl{Infinite Galois Theory.}  In
%\emph{Handbook of Algebra}, (M. Hazewinkel, Ed.) Vol. I, 269--319, Elsevier, Amsterdam (1996). 

%\bibitem{Kennan} S. Kennan. \textsl{Symmetries of tropical moduli spaces of curves.}
%\texttt{arXiv:2004.10912v1} (2020).

%\bibitem{[K]} M.Korkmaz. Automorphisms of complexes of curves on punctured spheres 
%and on punctured tori. Topology Appl. {\bf 95} (1999), 85-111.

%\bibitem{LS} P. Lochak, L. Schneps. \textsl{A cohomological interpretation of the Grothendieck-Teichm{\"u}ller group.} 
%Invent. Math. {\bf 127} (1997), 571--600.

%\bibitem{LNS} P. Lochak, H. Nakamura, L. Schneps. \textsl{On a new version of the Grothendieck-Teichm{\"u}ller group.} 
%C. R. Acad. Sci. Paris, t. {\bf 325} (1997), 11--16.

%\bibitem{[L]} F. Luo. \textsl{Automorphisms of the complex of curves.}
%Topology {\bf 39} (2000), 283-298.


%\bibitem{Margalit} D. Margalit. \textsl{Automorphisms of the pants complex.} Duke Math. J. {\bf 121} (2004), 457--479.

%\bibitem{MN}A. Minamide, H. Nakamura. \textsl{The automorphism groups of the profinite braid groups.}
%Amer. J. Math. {\bf 144} (5) (2022), 1159--1176.

%\bibitem{MZ} A. Minasyan, P. Zalesskii. \textsl{Virtually compact special hyperbolic groups are
%conjugacy separable}. Comment. Math. Helv. {\bf 91} (2016), 609--627.

%\bibitem{[Mo]} S. Mochizuki. \textsl{The local pro-$p$ anabelian geometry of curves.}
%Invent. Math. {\bf 138} (1999), 319--423.

%\bibitem{MTopics}S. Mochizuki. \textsl{Topics Surrounding the Anabelian Geometry of Hyperbolic Curves.}
%In: \emph{Galois groups and fundamental groups}. Vol. {\bf 41}. Math. Sci. Res. Inst. Publ. Cambridge Univ. Press, Cambridge (2003), 119--165.

%\bibitem{Moch1} S. Mochizuki. \textsl{Topics in Absolute Anabelian Geometry I: Generalities.}
%J. Math. Sci. Univ. Tokyo {\bf 19} (2012), 139--242.

%\bibitem{Moch2} S. Mochizuki. \textsl{Topics in Absolute Anabelian Geometry II:
%Decomposition Groups and Endomorphisms.}
%J. Math. Sci. Univ. Tokyo {\bf 20} (2013), 171--269.

%\bibitem{Moch3} S. Mochizuki. \textsl{Topics in Absolute Anabelian Geometry III:
%Global Reconstruction Algorithms.} J. Math. Sci. Univ. Tokyo {\bf 22} (2015), 939--1156.

%\bibitem{[McC]} J.D. McCarthy. \textsl{Automorphisms of surface mapping class groups:
%A recent theorem of N.Ivanov.} Invent. Math. {\bf 84} (1986), 49--71.

%\bibitem{[MNT]} S.Mochizuki, H.Nakamura, A.Tamagawa. The Grothendieck conjecture on the
%fundamental group of algebraic curves. Sugaku Expositions {\bf 14} (2001), 31-53.

%\bibitem{N} H. Nakamura. \textsl{Galois rigidity of pure sphere braid groups and profinite calculus.}
%J. Math. Sci. Univ. Tokyo {\bf 1} (1994), 72--136.

%\bibitem{[No]} B.Noohi. Fundamental groups of algebraic stacks.
%J. Inst. Math. Jussieu {\bf 3} (2004), 69-103.

%\bibitem {NS} H. Nakamura, L. Schneps. \textsl{On a subgroup of the Grothendieck-Teichm{\"u}ller group
%acting on the tower of profinite Teichm{\"u}ller modular groups.}
%Invent. Math. {\bf 141} (2000), 503--560.

%\bibitem{Pop1} F. Pop. \textsl{ON I/OM}. \url{https://www.math.upenn.edu/~pop/Research/files-Res/IOM_29May14.pdf} (2014).

%\bibitem{Pop2} F. Pop. \textsl{Finite tripod variants of I/OM}. Invent. Math. {\bf 216} (2019), 745--797.

%\bibitem {[Q1]} G.Quick.Group action on profinite spaces.
%J. Pure Appl. Algebra {\bf 215} (2011), 1024-1039.

%\bibitem {[Q2]} G.Quick. Some remarks on profinite completion of spaces.
%Advanced Studies in Pure Mathematics {\bf 63} (2012), 413-448.

%\bibitem{RZ}L. Ribes, P. Zalesskii. \textsl{Profinite groups. Second edition}. Erg. der Math.
%und ihrer Grenz. 3. Folge {\bf 40}, Springer-Verlag (2010).


%\bibitem{Ro}
%M. Romagny. \textsl{Group actions on stacks and applications}. Michigan Math. J. {\bf 53} (2005), 209--236. 

%\vskip .1cm
%\bibitem{[S]} K.J.Shackleton. Geodesic axes in the pants complex of the five-holed sphere.
%Topology and its Applications {\bf 158} (2011), 1640-1654.

%\bibitem{Ser1} 
%J.P.Serre, {\em Groupes de Galois sur $\Q$}, Sem. Bourbaki no. 689, 1987-1988. 
%\bibitem{Ser2}
%J.P.Serre, {\em Topics in Galois theory}, Research Notes in Math, AK Peters/CRC Press, 1992. 

%\bibitem{[ST]} M.Sa\"idi, A.Tamagawa. On the $Hom$-form of Grothendieck's birational conjecture
%in positive characteristic. Algebra \& Number Theory {\bf 5} (2011), 131-184.

%\bibitem{Sch}F. Schaffhauser. \textsl{Lectures on Klein Surfaces and Their Fundamental Group}. In: L. Alvarez Consul, J. Andersen,
%I. Mundet i Riera (eds) \emph{Geometry and Quantization of Moduli Spaces}. Advanced Courses in Mathematics - CRM Barcelona. 
%Birkh\"auser (2016), 67--108.

%\bibitem{Schneps} L. Schneps. \textsl{The Grothendieck-Teichm\"uller group $\GT$: a survey.} 
%In {\it Geometric Galois Actions I}, edited by L. Schneps and P. Lochak, Lon. Math. Soc. LNS {\bf 242} (1997), 183-203.

%\bibitem{ScSt} A. Schmidt, J. Stix. \textsl{Anabelian geometry with \'etale homotopy types.} Ann. of Math. (2), Vol. {\bf 184}, No. 3
%(2016), 817--868.

%\bibitem{Segal} I.E. Segal. \textsl{The automorphisms of the symmetric group}. Bull. Amer. Math. Soc. {\bf 46} (6) (1940), 565--565.

%\bibitem{stacks-project} The Stacks Project Authors. \textsl{The Stacks Project.} \url{https://stacks.math.columbia.edu} (2020).

%\bibitem{Vaintrob}D. Vaintrob.  \textsl{Moduli of framed formal curves.} 
%\url{https://arxiv.org/abs/1910.11550} (2019).

%\bibitem{Wolpert}S. A. Wolpert.  \textsl{The Weil-Petersson metric geometry.} \url{https://arxiv.org/pdf/0801.0175.pdf} (2007).


\end{thebibliography}

\end{document}










Let us then see what happens when we fill in a set of punctures which are unramified for the given cover. Let ${\mathcal Q}$ be a set of points of 
$S_{/G}^\circ$ and $\wt{\cQ}:=p_G^{-1}(\cQ)$. Let $S^{\wt{\cQ}}:=S\ssm\wt{\cQ}$ and $S_{/G}^\cQ:=S_{/G}^\circ\ssm \cQ$. Put then $q:=\sharp \cQ$.

For a given surface $S$, let us denote by $S[n]$ the unordered configuration space of $n$ points on the surface $S$.  
The set of points $\cQ$ determines a point on $S_{/G}^\cQ[q]$, which we also denote by $\cQ$. The fundamental group $\pi_1(S_{/G}^\cQ[q],\cQ)$
identifies with a subgroup of the mapping class group $\G(S_{/G}^\cQ)$. 
Let then $\pi_1(S_{/G}^\cQ[q],\cQ)_K:=\pi_1(S_{/G}^\cQ[q],\cQ)\cap\G(S_{/G}^\cQ)_K$. We have:

\begin{proposition}\label{unforg}The normalizer $N_{\G(S^{\wt{\cQ}})}(G)$ of the finite group $G$ in $\G(S^{\wt{\cQ}})$ fits in 
the short exact sequence:
\[1\to\pi_1(S_{/G}^\cQ[q],\cQ)_K\to N_{\G(S^{\wt{\cQ}})}(G)\to N_{\G(S)}(G)\to 1.\]
\end{proposition}

\begin{proof}The normalizers $N_{\G(S^{\wt{\cQ}})}(G)$ and $N_{\G(S)}(G)$ fit in the exact commutative diagram:
\[\begin{array}{ccccc}
&1&&1&\\
&\da&&\da&\\
\hspace{0.85cm}1\to&\ker\psi&\sr{\sim}{\to}&\pi_1(S_{/G}^\cQ[q],\cQ)_K&\to 1\\
\hspace{0.4cm}\da&\da&&\da&\\
1\to G\to& N_{\G(S^{\wt{\cQ}})}(G)&\to&\G(S_{/G}^\cQ)_K&\to 1\\
\hspace{0.4cm}\parallel&\da^\psi&&\da&\\
1\to G\to &N_{\G(S)}(G)&\to&\G(S_{/G}^\circ)_K&\to 1\\
\hspace{0.4cm}\da&\da&&\da&\\
\hspace{0.4cm}1&1&&\,1.&
\end{array}\]
From which, the conclusion follows.
\end{proof}
